
\begin{figure}[t]
\centering
  \includegraphics[width=\linewidth]{figures/mdr_combined.pdf}
  \caption{Layer-wise MDR analysis on LLaVA-1.6-7B. \textbf{Left:} Non-Conflict and Conflict-Correct samples maintain positive MDR, while Conflict-Incorrect samples reverse from positive to negative in late layers. \textbf{Right:} Text-ward shifts ($\Delta$MDR$<$0) dominate failures (85\%), vision-ward shifts ($\Delta$MDR$>$0) dominate correct predictions (89--91\%). Shading: $\pm$1 std.}
  \label{fig:mdr_trajectory}
\end{figure}

\section{Diagnosing Late-Layer Textual Override}
\label{sec:analysis}

To understand why models favor text over vision under conflict, we need to trace how their predictions evolve during the forward pass. Does textual bias arise because visual information is never properly encoded? Or is it encoded but subsequently discarded? This section presents a layer-wise analysis that reveals a surprising answer.

\subsection{Experimental Framework}
\label{sec:setup}

\textbf{Task and data.} We study visual question answering where inputs consist of image $I$, context $C$, and question $Q$. Let $w_{\text{vis}}$ and $w_{\text{text}}$ denote answers supported by the image and context respectively. A \emph{knowledge conflict} occurs when $w_{\text{vis}} \neq w_{\text{text}}$. We categorize samples into three groups based on context type and model output. \textbf{Non-Conflict}: the context is consistent with the image and the model answers correctly. \textbf{Conflict-Correct}: context contradicts image but the model follows visual evidence. \textbf{Conflict-Incorrect}: the context contradicts the image and the
model follows textual claims.

Existing conflict benchmarks lack the $(w_{\text{vis}}, w_{\text{text}})$ annotations required for layer-wise analysis. To address this, we construct \textbf{Conflict-VQA}: 5,969 samples with paired consistent and conflicting contexts. We source images from VrR-VG~\cite{VrR-VG2019} and questions from TDIUC~\cite{TD-IUC2017}, covering six question types: object
presence, color, attribute, counting, spatial reasoning, and activity recognition. All contexts
are generated using GPT-4o and manually verified. Construction details are provided in Appendix A. For mechanistic analysis, we filter each model's competent subset, defined as samples answered correctly under consistent context across five trials. This filtering ensures that observed failures under conflict reflect conflict resolution process rather than basic visual perception errors.

\textbf{Layer-wise projection.} Following early-exit methods~\cite{depth_adaptive2020,confident_adaptive2022}, we project hidden states at each layer $l$ to the output vocabulary via the language model head $W_{\text{head}}$: $P^{(l)}_t=\text{softmax}(W_{\text{head}}h^{(l)}_t)$, where $h^{(l)}_t$ is the hidden state at the answer token position and $P^{(l)}_t(w)$ is the scalar probability assigned to vocabulary token $w$. This lets us observe the model's ``preference'' at each depth.


\begin{figure}[t]
\centering
  \includegraphics[width=\linewidth]{figures/compare.pdf}
  \caption{JSD and MDR trajectories on LLaVA-1.6-7B. Both categories show JSD spikes at similar depths, but MDR moves in opposite directions: toward vision for Non-Conflict, toward text for Conflict-Incorrect.}
  \label{fig:jsd_mdr}
\end{figure}


\subsection{Visual Predictions Emerge Early but Fade in Late Layers}
\label{sec:discovery1}


When models fail under conflict, is visual information absent from their representations, or present but ultimately ignored?

To answer this, we introduce \textbf{Modal Dominance Ratio (MDR)}, quantifying modal preference at each layer:
\begin{equation}
\text{MDR}^{(l)} = \log \frac{P^{(l)}_t(w_{\text{vis}})}{P^{(l)}_t(w_{\text{text}})}
\end{equation}
Positive MDR indicates preference for the visual answer; negative indicates textual preference.

Figure~\ref{fig:mdr_trajectory} shows MDR trajectories averaged across sample categories.
Non-Conflict and Conflict-Correct samples maintain positive MDR throughout all layers, reflecting consistent visual preference. Conflict-Incorrect samples tell a different story: \textbf{MDR is positive in early and middle layers, then reverses to negative in late layers.}

This pattern indicates that in failure cases, models assign higher probability to the visual answer in intermediate layers before shifting toward the textual answer. In other words, the visual information is present in intermediate representations but does not survive to the final output. We term this phenomenon \emph{late-layer textual override}. This observation suggests that the failure mechanism lies in late-layer processing rather than in early-layer encoding.


\subsection{Transition Direction Distinguishes Success from Failure}
\label{sec:discovery2}


The late-layer reversal in Conflict-Incorrect samples might suggest a straightforward fix: simply use intermediate-layer predictions. But prediction transitions also occur when models answer correctly. In these cases, late-layer processing appears to refine the answer rather than override it. To intervene selectively, we need a way to tell these two situations apart.


\begin{figure}[t]
\centering
    \includegraphics[width=\linewidth]{figures/detection.pdf}
    \caption{Detection signal distribution on LLaVA-1.6-7B. Conflict-Incorrect samples (red) exhibit high $D_{\mathrm{conf}}$ but low $\rho$. Non-Conflict (blue) and Conflict-Correct (green) maintain high $\rho$.}
    \label{fig:detection_signals}
\end{figure}


\textbf{Detecting transitions.} Layer-wise predictions show high entropy in early layers and lower entropy in deeper layers as predictions converge.
We define the \emph{stability onset} $L_{\text{start}}$ as the layer with the largest entropy drop:
\begin{equation}
L_{\text{start}} = \arg\max_l [H_t^{(l-1)} - H_t^{(l)}],
\end{equation}
where $H_t^{(l)} = -\sum_w P^{(l)}_t(w) \log P^{(l)}_t(w)$.
Across our models, $L_{\text{start}}$ falls between 25\% and 40\% of model depth.


After stability onset, we detect abrupt distributional shifts using Jensen-Shannon divergence between adjacent layers: \begin{equation}
\text{JSD}^{(l)} = \frac{1}{2}\left[D_{\text{KL}}(P^{(l)}_t \| M) + D_{\text{KL}}(P^{(l+1)}_t \| M)\right],
\end{equation}
where $M = \frac{1}{2}(P^{(l)}_t + P^{(l+1)}_t)$. The transition layer is: $L_{\text{trans}} = \arg\max_{l \in [L_{\text{start}}, N-1]} \text{JSD}^{(l)}$.

\textbf{Transition direction distinguishes success from failure.} Figure~\ref{fig:jsd_mdr} shows that both Non-Conflict and Conflict-Incorrect samples have JSD spikes at similar depths, so transitions alone do not distinguish the two groups. What does differ is the \emph{direction} of the shift. We quantify this via MDR change across the transition:
\begin{equation}
\Delta\text{MDR} = \overline{\text{MDR}}_{\text{after}} - \overline{\text{MDR}}_{\text{before}}
\end{equation}
where bars denote averages over pre/post-transition layers.


\smallskip\noindent\textbf{Finding:} In 85\% of Conflict-Incorrect cases, $\Delta\text{MDR}$ is negative, indicating a shift toward text. The pattern reverses for correct predictions: 89\% of Conflict-Correct and 91\% of Non-Conflict samples show positive $\Delta\text{MDR}$, shifting toward vision or holding steady. Transitions occur in all groups, but only the direction correlates with correctness.
